\documentclass[10pt,letterpaper]{article}
\usepackage[T1]{fontenc}
\usepackage[utf8]{inputenc}
\usepackage{lmodern,microtype}
\usepackage[margin=0.85in]{geometry}
\usepackage{amsmath,amssymb,amsthm,mathtools}
\usepackage{booktabs,longtable,array}
\usepackage{xcolor,fancyhdr,enumitem}
\usepackage{xurl}
\usepackage[colorlinks=true,linkcolor=blue!45!black,citecolor=blue!45!black,urlcolor=blue!45!black]{hyperref}
\hypersetup{pdftitle={A direct quartic for Con(ZFC) in 234 integer variables},pdfauthor={AI-assisted research prepared for Vladimir Reshetnikov and ProveIt},pdfsubject={Exact quadratic elimination and cofinal arithmetic witnesses}}
\setlength{\headheight}{14pt}
\pagestyle{fancy}\fancyhf{}
\fancyhead[L]{\small A direct quartic in 234 integer variables}
\fancyhead[R]{\small ProveIt research contribution}
\fancyfoot[C]{\thepage}
\setlength{\parindent}{0pt}\setlength{\parskip}{0.4em}
\setlist{nosep,leftmargin=*}
\allowdisplaybreaks[3]\emergencystretch=2em
\newtheorem{theorem}{Theorem}[section]
\newtheorem{lemma}[theorem]{Lemma}
\newtheorem{proposition}[theorem]{Proposition}
\newtheorem{corollary}[theorem]{Corollary}
\theoremstyle{definition}\newtheorem{remark}[theorem]{Remark}
\newcommand{\Z}{\mathbb Z}\newcommand{\Q}{\mathbb Q}\newcommand{\N}{\mathbb N}
\newcommand{\Con}{\operatorname{Con}}
\newcommand{\Par}{P_{\mathrm{240}}}
\newcommand{\New}{P_{\mathrm{234}}}
\newcommand{\BP}{P_{\mathrm{235},\beta}}
\newcommand{\SP}{P_{\mathrm{236},\beta}}
\title{\textbf{A direct quartic for $\Con(\mathrm{ZFC})$\\in 234 integer variables}\\[0.5em]
\large Exact elimination along common quadratic directions\\and cofinal quadratic spacing}
\author{AI-assisted mathematical research\\Prepared for Vladimir Reshetnikov and the ProveIt project}
\date{25 October 2026}
\begin{document}
\maketitle\thispagestyle{empty}
\begin{abstract}
We give an explicit polynomial $\New\in\Z[x_0,\ldots,x_{233}]$ whose
integer unsolvability is provably equivalent in ZFC to $\Con(\mathrm{ZFC})$,
by composing a new exact reduction with the pinned ProveIt numerical-calculus
calibration. It is a sum of 171 explicitly listed quadratic squares, has
degree four, 19,050 collected monomials, and coefficient height 4,480.
Its actual parent has 240 variables, 177 residuals, 21,482 monomials, and
height 6,144. Thus all four size measures improve while degree stays four.

The main reduction identifies six integral directions annihilated by every
quadratic Hessian. A unit linear slope permits deletion of one variable and
one residual, with an explicit integer polynomial inverse on zero sets.
An independently implemented checker verifies all six transformations and
every coefficient of the final polynomial. A separate arithmetic change
replaces two beta-spacing equations by $y=F^2+u^2+v^2$. An elementary
prime-power and CRT argument makes this family cofinal among suitable
factorial multiples. This produces factorial-free companions in 235 and
236 integer variables.

The contribution is a concrete reduction of the direct small-coefficient
branch, not a global minimum or a new universal-variable bound. The algebra
and cofinality arguments are ordinary unbounded proofs; finite coefficient
verification and the inherited logical calibration have distinct scopes.
No end-to-end Lean or Rocq proof is claimed.
\end{abstract}
\vfill
\begin{center}\small
Comparison snapshot: \texttt{7a8da414376a80b90ecf799f37ceb26fb9ac674e}.\\
All final variables range over the unrestricted integers.\\
Every principal residual is printed in Appendix~\ref{app:residuals}.
\end{center}
\newpage\tableofcontents\newpage

\section{The explicit result}

The principal mathematical object is the fixed ordinary polynomial
\begin{equation}\label{eq:main}
 \boxed{\qquad \New(x_0,\ldots,x_{233})=
            \sum_{i=0}^{170}\mathcal Q_i(x_0,\ldots,x_{233})^2.\qquad}
\end{equation}
Appendix~\ref{app:residuals} specifies every coefficient of every
$\mathcal Q_i$. The file \nolinkurl{equations/P234/equation.json} is exactly
the same definition in sparse machine-readable form. Its neighboring
\nolinkurl{expanded_polynomial.json} lists all 19,050 collected terms of
$\New$ directly. Neither definition contains a variable exponent, binomial
coefficient, proof-checking oracle, unspecified universal polynomial, or
uncomputed program-dependent integer.

For orientation, the first six residuals are
\begin{align*}
\mathcal Q_0&=x_4-x_3^2,\\
\mathcal Q_1&=x_0-x_0x_2-x_2+x_9,\\
\mathcal Q_2&=-x_2x_7-x_3+x_8x_9,\\
\mathcal Q_3&=-1+x_1-x_8x_9,\\
\mathcal Q_4&=1+x_4-x_5x_{20}+x_{18},\\
\mathcal Q_5&=-1-x_{10}^2-2x_{10}x_{11}-x_{11}^2+x_{21}.
\end{align*}
The other 165 residuals are part of the definition, not omitted hypotheses.
Every residual coefficient has absolute value at most 64. After squaring
and collecting, the coefficient of $x_0^4$ is 13, so the degree is exactly
four. The largest absolute coefficient is attained, for example, at
$x_{11}^2x_{22}^2$, whose coefficient is 4,480.

\begin{theorem}[Principal consistency equation]\label{thm:main}
For the polynomial \eqref{eq:main}, the new algebraic proof composed with
the pinned ProveIt calibration establishes
\[
 \mathrm{ZFC}\vdash
 \left(\neg\exists\boldsymbol x\in\Z^{234}\quad\New(\boldsymbol x)=0\right)
 \ \longleftrightarrow\ \Con(\mathrm{ZFC}).
\]
More specifically, ZFC proves an explicit integer zero-set bijection
between the new residual system and the pinned 240-variable parent.
No consistency assumption is used to establish this reduction.
\end{theorem}

The exact new result proved independently of the inherited semantics is
\begin{equation}\label{eq:algebra-calibration}
 \mathrm{ZFC}\vdash
 \bigl(\exists\boldsymbol x\in\Z^{234}\;\New(\boldsymbol x)=0\bigr)
 \longleftrightarrow
 \bigl(\exists\boldsymbol z\in\Z^{240}\;\Par(\boldsymbol z)=0\bigr).
\end{equation}
The parent theorem identifies the latter existence statement with finite
syntactic ZFC refutability. Section~\ref{sec:logical} states that dependency
precisely; it is not silently reclassified as a new formal verification.

\subsection{Measured comparison and the objective}

\begin{center}\small
\begin{tabular}{lrrrrr}
\toprule
Equation & Variables & Residuals & Degree & Monomials & Height\\
\midrule
Shared-error parent $\Par$ &240&177&4&21,482&6,144\\
\textbf{Principal $\New$} &\textbf{234}&\textbf{171}&\textbf{4}&\textbf{19,050}&\textbf{4,480}\\
Factorial-free parent &242&178&4&5,018&7,968\\
Factorial-free $\BP$ &235&170&4&18,550&4,576\\
Smaller-support $\SP$ &236&171&4&4,947&16,866\\
\bottomrule
\end{tabular}
\end{center}

Height means the largest absolute coefficient of the fully collected
polynomial, rather than of an uncollected expression or of one residual.
Monomials are counted after all cancellations. The three delivered outputs
have different advantages: $\New$ has the smallest dimension and height;
$\BP$ uses the factorial-free arithmetic source and has one fewer residual;
$\SP$ has substantially smaller support.

Relative to its actual parent, $\New$ removes six variables and six
residuals, 2,432 collected monomials, and 1,664 from the coefficient height.
The reductions are 2.5\%, approximately 11.32\%, and approximately 27.08\%
for variables, monomials, and height, respectively. This is a concrete
improvement of that branch, not a major new universal bound.

The inspected repository already contains nine-positive-variable and
other low-variable constructions, including quartics in 37 unrestricted
integers on a different branch~\cite{overview}. Consequently 234 is not a
project-wide variable record. The present objective is a direct numerical
proof-calculus interpretation, degree four, small literal coefficients,
and an auditable finite definition. No scalar measure called
``simplest'' orders all these resources. The distinction between degree,
variable count and explicit representation is also central in current
work on universal Diophantine bounds~\cite{bayer,afp}.

\section{What the parent certifies}\label{sec:logical}

The comparison source is the unchanged incoming archive
\emph{Shared Factorial Error}, at the pinned commit~\cite{shared,repo}.
Its selected file is
\nolinkurl{equations/equation240/equation.json}. The archive hash is
recorded in \nolinkurl{SOURCES.json}, and the exact member bytes are checked
against the bundled input on every release verification.

The parent uses the numerical calculus $C_{222}$, with truth code 0 and
falsity code 1. It represents a finite multiset of single-conclusion Horn
rule instances together with a balance condition. Bounded row certificates
and a Chinese remainder construction enforce the local arithmetic. A
shared Pell construction represents the required powers by ordinary
polynomial equations. The parent does not use a simulation of K215 or a
specialization of Jones's universal polynomial in its emitted definition.

The logical endpoints inherited from ProveIt are
\begin{center}\small
\nolinkurl{KEndpoint.c222_iff_zf_refutation},\\
\nolinkurl{ConstructibleChoice.zf_refutation_iff_zfc_refutation}.
\end{center}
They identify finite $C_{222}$ refutability with syntactic ZF refutability
and then syntactic ZFC refutability. The ZF/ZFC step uses the established
constructible-universe equiconsistency argument. The retained source
archives identify the numerical rules and ordinary arithmetic-to-calculus
proofs as well as these logical interfaces.

A useful reason that the multiset representation is sound is the following
finite closure argument. Pad unused premise slots by derivable truth.
Suppose at every positive label the number of conclusions minus premises
is zero, except at the target label where it is one. If the target were
outside the finite Horn closure of truth, every rule with an outside
conclusion would have at least one outside premise. Summing conclusions
minus premises over the outside labels would give a number at most zero,
contradicting one. Thus circular collections cannot manufacture a false
refutation merely by balancing multiplicities.

In the shared-error source, the arithmetic quantity $F$ is authenticated as
$t!$ by a nearest-integer fixed point. The normalized common error supplies
\[
 \left|F-\frac{(F^2+1)^t}{\binom{F^2+t}{t}}\right|<\frac12.
\]
For $t\ge16$ and the signs recovered in the parent, its unbounded proof
forces $F=t!$. A shared norm simultaneously controls this error and eleven
digit remainders. The signed-base power lemma is used before signs of the
radix and binomial output are known. These ordering issues matter: a
positivity conclusion obtained from a decoded digit cannot be assumed in
the estimate that first decodes that digit.

The independent ordinary review in
\nolinkurl{proofs/parent_arithmetic_review.md} checks the changed signed
bootstrap, the separate unit-base and $F=0$ cases, the nearest-factorial
estimate, and the compatible witness construction. It found no material
gap in those arguments under their retained Pell and calibration
interfaces. The original release verifier was also rerun successfully,
including all 242 sealed files, its earlier source transformations, the
240-variable profile, and its printed equation.

These checks do not constitute a fresh build of the entire transitive
Lean or Rocq development. The new contribution proves
\eqref{eq:algebra-calibration} by elementary algebra, and then composes
with the identified ordinary parent theorem. This explicit dependency is
the mathematical scope of Theorem~\ref{thm:main}.

\section{Elimination along a common quadratic direction}

Write an integer quadratic polynomial as
\[
 Q_i(\boldsymbol x)=q_i(\boldsymbol x)+\ell_i(\boldsymbol x)+c_i,
\]
where $q_i$ is homogeneous of degree two, $\ell_i$ is linear, and $c_i$
is constant. Let $H_i$ be its integer Hessian: the diagonal entry for a
term $a x_j^2$ is $2a$, and a mixed term $b x_jx_k$ contributes $b$ at
each of the two symmetric off-diagonal positions.

\begin{lemma}[Constant directional slope]\label{lem:direction}
If $\boldsymbol v\in\Z^n$ satisfies $H_i\boldsymbol v=0$, then for all
$\boldsymbol y\in\Z^n$ and $s\in\Z$,
\[
 Q_i(\boldsymbol y+s\boldsymbol v)
      =Q_i(\boldsymbol y)+s\,\ell_i(\boldsymbol v).
\]
\end{lemma}
\begin{proof}
Expanding the homogeneous quadratic part gives the cross term
$s\boldsymbol y^TH_i\boldsymbol v$ and the pure term
$s^2q_i(\boldsymbol v)$. The former vanishes. Since
$2q_i(\boldsymbol v)=\boldsymbol v^TH_i\boldsymbol v=0$, the latter
vanishes as well. The linear and constant parts give the stated formula.
\end{proof}

\begin{theorem}[Integral quadratic projection]\label{thm:projection}
Let $Q_0,\ldots,Q_{m-1}\in\Z[x_0,\ldots,x_{n-1}]$ have degree at most
two. Suppose there are $\boldsymbol v\in\Z^n$ and indices $p,r$ such that
\[
 H_i\boldsymbol v=0\quad\text{for every }i,\qquad v_p=1,\qquad
 a_r:=\ell_r(\boldsymbol v)\in\{-1,1\}.
\]
Set $a_i=\ell_i(\boldsymbol v)$. Regard $\boldsymbol y$ as a vector with
$y_p=0$ and $n-1$ free integer coordinates, and define
\begin{equation}\label{eq:projected}
 R_i(\boldsymbol y)=Q_i(\boldsymbol y)-a_i a_rQ_r(\boldsymbol y),
 \qquad i\ne r.
\end{equation}
These are integer quadratics. Their common integer zero set is in explicit
bijection with the common integer zero set of the original $Q_i$.
\end{theorem}
\begin{proof}
Given a new zero, put
\begin{equation}\label{eq:lift}
 s=-a_rQ_r(\boldsymbol y),\qquad
 \boldsymbol x=\boldsymbol y+s\boldsymbol v.
\end{equation}
Every coordinate is an integer. Lemma~\ref{lem:direction} gives
$Q_r(\boldsymbol x)=Q_r(\boldsymbol y)+sa_r=0$, because $a_r^2=1$.
For every other row it gives
$Q_i(\boldsymbol x)=R_i(\boldsymbol y)=0$.

Conversely, given an old zero, let $s=x_p$ and
$\boldsymbol y=\boldsymbol x-x_p\boldsymbol v$. Then $y_p=0$.
The same identity gives $Q_r(\boldsymbol y)=-x_pa_r$, so the scalar
reconstructed by \eqref{eq:lift} is exactly $x_p$. For every $i\ne r$,
\[
 R_i(\boldsymbol y)=-x_pa_i-a_i a_r(-x_pa_r)=0.
\]
The two constructions are inverse. Formula~\eqref{eq:projected} is a
constant integer row combination restricted to $y_p=0$, hence has degree
at most two. There is no division or residual congruence condition.
\end{proof}

\begin{corollary}
Replacing $\sum_iQ_i^2=0$ by $\sum_{i\ne r}R_i^2=0$ removes one integer
unknown and one quadratic residual without increasing degree above four.
The equivalence is provable by finite arithmetic in ZFC.
\end{corollary}
\begin{proof}
A sum of integer squares is zero precisely when each summand is zero.
Apply Theorem~\ref{thm:projection}.
\end{proof}

\subsection{A small example}

Consider the residuals
\[
 u-a^2,\qquad v-b^2,\qquad w-(u+v)c.
\]
The direction that increases $u$ and decreases $v$ annihilates all three
quadratic Hessians. With $s=u+v$, the system becomes
\[
 s-a^2-b^2=0,\qquad w-sc=0.
\]
It has one fewer integer variable and one fewer residual. A zero of the
new system restores $u=a^2$, $v=b^2$, and the inverse uses $s=u+v$.
A procedure that only substitutes a single existing variable will initially
see cubic terms in $uc$ or $vc$. Changing to the combination $u+v$ exposes
the quadratic-preserving reduction.

\subsection{Integrality is an essential hypothesis}

A primitive integral direction alone does not justify division by its
slope. If all slopes have a common factor greater than one, solving for
the eliminated coordinate can leave a divisibility condition. Our
certificates check $v_p=1$ and a literal slope $\pm1$ at every step.
More general primitive directions and slope vectors can be handled by
unimodular completions and Bezout combinations, but that generalization
is not needed for any claimed count in this release.

The proof also uses the unrestricted integer domain. An invertible integer
change of coordinates need not preserve positivity or nonnegativity.
The present reduction is not automatically a reduction of the repository's
nine-positive-variable equations.

\section{The six eliminations and the coefficient reduction}

The discovery program stacks the sparse Hessian rows and computes their
common kernel by exact rational Gaussian elimination. It clears
denominators, makes each vector primitive, and considers coordinates with
unit entries. It then checks every Hessian product and every candidate
unit slope exactly. The selected coordinate change and unit substitution
are stored as a finite certificate.

For the principal parent, the six deleted register names, in order, are
\begin{center}\small
\begin{tabular}{crrrl}
\toprule
Step & Input variables & Output variables & Output residuals & Deleted register\\
\midrule
1&240&239&176&\texttt{add.71}\\
2&239&238&175&\texttt{mul.188}\\
3&238&237&174&\texttt{mul.303}\\
4&237&236&173&\texttt{sqd.11}\\
5&236&235&172&\texttt{add.277}\\
6&235&234&171&\texttt{add.112}\\
\bottomrule
\end{tabular}
\end{center}
These are historical wire names, not mathematical operations remaining in
the final polynomial. The coordinate indices, complete integer direction
vectors, selected residuals, unit slopes and restoration polynomials are
specified in \nolinkurl{stages/radical234/projection_certificate.json}.
Renumbering after each deletion is recorded explicitly.

Immediately after these six steps the quartic has 20,739 monomials and
height 18,542. We then apply 147 recorded elementary operations
\begin{equation}\label{eq:rowop}
 R_i\ \longleftarrow\ R_i+kR_j,
 \qquad i\ne j,\quad k\in\Z.
\end{equation}
Each is invertible by subtracting $kR_j$. Thus the common zero set is
unchanged, although the polynomial obtained by summing squares usually
changes. The operations give the final 19,050 monomials and height 4,480.
They are recorded in \nolinkurl{equations/P234/row_certificate.json}.

The row search chooses reductions of the squared Euclidean norm of the
coefficient vectors. Its heuristic has no role in the correctness proof.
The checker replays the frozen operations exactly and then replays them
backwards, recovering the original intermediate rows. This is the reason
that a changed optimization heuristic need not be trusted or rerun.

The independent projection checker uses \eqref{eq:projected} directly.
The constructor uses a shear followed by substitution. These two
implementations agree coefficient by coefficient on every intermediate
residual. Numerical evaluations at arbitrary small tuples supplement
these identities as regression controls; they are not substitutes for
the identities and are not assumed zeros of the consistency equation.

\section{A second contribution: cofinal quadratic spacing}\label{sec:cofinal}

The factorial-free parent~\cite{factorialfree} uses
\[
 t=16+h^2,\qquad F=[2(t+1)]^t,
\]
where the displayed power is represented by the inherited ordinary Pell
certificate. Its beta spacing is specified by two quadratic register rows
\begin{equation}\label{eq:oldbeta}
 m=(2w+1)^2,\qquad y=Fm.
\end{equation}
The spacing must be large enough to encode a finite accumulator path; in
the converse construction it is chosen as a multiple of $t!$. Authenticating
$F=t!$ is not part of this branch.

We replace \eqref{eq:oldbeta} by the single quadratic row
\begin{equation}\label{eq:newbeta}
 \boxed{\qquad y=F^2+u^2+v^2,\qquad u,v\in\Z.\qquad}
\end{equation}
The already stored value $F^2$ is reused. This uses the same two witness
slots, removes one residual, and implies $y\ge F^2>F+1$ once the existing
growth certificate has been interpreted. The following theorem supplies
exactly the completeness property needed.

\begin{lemma}[Two squares modulo an odd prime power]\label{lem:odd}
For every odd prime $p$ and $e\ge1$, there exist integers $a,b$ with
$a^2+b^2\equiv-1\pmod{p^e}$.
\end{lemma}
\begin{proof}
The squares in $\mathbb F_p$ and minus one minus those squares are two
sets of size $(p+1)/2$, so they intersect. A resulting solution modulo $p$
cannot have both coordinates zero. Suppose the first is nonzero; the other
case is symmetric. Given a solution modulo $p^j$, change $a$ to
$a+\delta p^j$, where
\[
 \delta\equiv
 -\frac{1+a^2+b^2}{p^j}(2a)^{-1}\pmod p.
\]
Expansion proves the congruence modulo $p^{j+1}$. The coordinate remains
nonzero modulo $p$, so induction completes the lift.
\end{proof}

\begin{theorem}[Cofinal factorial multiples]\label{thm:cofinal}
Let $t\ge1$, $F\ge2$, and assume $2^{v_2(t!)}\mid F^2$. For every
integer bound $C$, there are integers $u,v$ such that
\[
 y=F^2+u^2+v^2>C,\qquad t!\mid y.
\]
Both $u$ and $v$ can be chosen as multiples of $F$. If $F$ is even,
the resulting $y$ is even.
\end{theorem}
\begin{proof}
Write $t!=2^a m$ with $m$ odd. Lemma~\ref{lem:odd} and the Chinese
remainder theorem give $m\mid1+A^2+B^2$. This also holds when $m=1$,
with the congruence vacuous. For positive integers $k$, put
\[
 u_k=F(A+km),\qquad v_k=FB,\qquad
 y_k=F^2\bigl(1+(A+km)^2+B^2\bigr).
\]
The odd part $m$ divides $y_k$, and $2^a$ divides $F^2$. Hence $t!$ divides
$y_k$. The family is unbounded, so some $y_k$ exceeds $C$. Its parity
when $F$ is even is immediate.
\end{proof}

For the actual $F=[2(t+1)]^t$, we have $v_2(F)\ge t>v_2(t!)$, so the
hypothesis holds with a large margin. The theorem is elementary and
unbounded. The delivered program constructs 195 exact examples for
$t=16,\ldots,80$ and three offsets each; those examples are checks of the
constructor, not the proof of cofinality.

\subsection{Soundness in the required order}

The known positive horizon and positive base first recover
$F=[2(t+1)]^t$ from the unchanged power certificate. Only then does
\eqref{eq:newbeta} imply $y\ge F^2>F+1$. The retained radix is an odd
square $q=(2z+1)^2$, hence at least one. Its numerator power requests at
base $q+1\ge2$ make the nine actual CRT representatives nonnegative;
therefore their offset sum $H$ is at least two.

The retained index row has the form
$B\iota=C_0+F^2+1$, with a nonnegative representative $C_0$, and excludes
$B=0$. Thus $|BH|\ge2$. The signed-power lemma gives
$x\ge0$ and $q=(BH)^x$. The endpoint
\[
 x=1+[1+(t+1)y]u_t
\]
excludes $x=0$, since its bracket is greater than one and cannot divide
$-1$. In the other endpoint,
\[
 x=F+1+(y+1)u_0,
\]
a negative integer $u_0$ would force $x\le F-y<0$. Hence
\[
 u_0\ge0,\qquad x\ge F+1,\qquad
 u_t\ge1,\qquad x\ge2+(t+1)y\ge2+(t+1)F^2.
\]
These bounds strengthen the parent bounds. Its centered-remainder proof
can therefore recover the positive binomial value
$B=\binom{F^2+t}{t}$, and its generalized CRT theorem recovers every
bounded local rule. No sign of that binomial output is assumed before
its digit is authenticated.

The initial accumulator $F+1$ is now a proper remainder modulo $y+1$
immediately. The old boundary case $y=F$, which required a zero-path
contradiction, is impossible. The same positive-weight product and finite
Horn-closure argument yield a $C_{222}$ refutation.

\subsection{Completeness and the parity choices}

Start with a finite $C_{222}$ refutation. Choose a sufficiently large
permitted horizon $t=16+h^2$, form $F$, and construct its finite positive
accumulator path before choosing $y$. Theorem~\ref{thm:cofinal} supplies
an even multiple of $t!$ exceeding every accumulator and $F+1$.
The moduli $1+(i+1)y$ for $0\le i\le t$ are pairwise coprime: a common
prime divides the index difference, hence is at most $t$ and divides $y$,
contrary to its dividing $1+(i+1)y$. They are all odd.

CRT gives the required beta code $x$. Adding their odd product changes
parity, so $x$ can be chosen even and arbitrarily large. Define the same
$B$ and $N=B^x$ as in the parent and use its generalized binomial CRT at
every prime of $B$. The value $B$ is odd because
$v_2(F^2)>\lfloor\log_2t\rfloor$: each $F^2+i$ has the same two-adic
valuation as $i$, so those valuations cancel in the binomial quotient.
The precision modulus is therefore odd and its representative classes
have free parity. The parent's choices of a dominating even anchor and
odd $H$ remain available. Then $q=(BH)^x$ is an odd square.

The other power, remainder, and norm witnesses are supplied by the
unchanged constructive arguments. Thus the new spacing preserves the
existence equivalence with a finite refutation. It generally does
\emph{not} extend an individual zero to the old block
\eqref{eq:oldbeta}: $y/F$ need not be an odd square. We use an existential
soundness/completeness proof for this step, rather than claiming a
coordinate bijection.

\section{The two factorial-free companions}

In the parent's final 242-variable chart, the actual spacing is $x_2-1$,
$F=x_3$, and its existing square is $x_4$. The two private witness slots
are 240 and 241. The literal source change is
\begin{align*}
 -1+x_2-x_3x_{240}&\quad\rightsquigarrow\quad
 -1+x_2-x_4-x_{240}^2-x_{241}^2,\\
 -1+x_{240}-4x_{241}-4x_{241}^2&\quad\text{is deleted}.
\end{align*}
The row $x_4-x_3^2$ is retained. The checker verifies that neither private
witness occurs anywhere else. The old projection and row certificates
also leave this block untouched; the arithmetic review records this
transport explicitly.

The new source has 242 variables and 177 residuals. Six applications of
Theorem~\ref{thm:projection} give 236 variables and 171 residuals. One
further elementary row operation gives $\SP$, with 4,947 monomials and
height 16,866. The definition is
\[
 \SP=\sum_{i=0}^{170}S_i^2,
\]
where every $S_i$ is listed literally in
\nolinkurl{equations/P236_beta/equation.json}.

For the 235-variable companion, an affine residual in the 236-variable
intermediate chart permits the following exact substitution:
\begin{align*}
 x_{96}={}&x_9+x_{10}+4x_{20}+2x_{23}+2x_{27}+2x_{29}+2x_{31}\\
 &+2x_{33}+4x_{35}+2x_{36}+2x_{39}+2x_{42}+4x_{61}+2x_{62}.
\end{align*}
Its unit coefficient makes the restoration integral and unique.
Deleting that row and substituting everywhere leaves 235 variables and
170 residuals, still quadratic. After 433 reversible integer row
operations, the resulting quartic $\BP$ has 18,550 monomials and height
4,576. All its coefficients appear in
\nolinkurl{equations/P235_beta/equation.json} and its expanded neighbor.

\begin{corollary}
The ordinary new proofs, composed with the factorial-free parent's
calibration, give in ZFC both equivalences
\[
 \neg\exists\boldsymbol x\in\Z^{235}\;\BP(\boldsymbol x)=0
       \ \longleftrightarrow\ \Con(\mathrm{ZFC}),
 \qquad
 \neg\exists\boldsymbol x\in\Z^{236}\;\SP(\boldsymbol x)=0
       \ \longleftrightarrow\ \Con(\mathrm{ZFC}).
\]
\end{corollary}

\section{Verification and its limits}

\subsection{Exact finite statements that were checked}

The independent checker is a standard-library Python program and imports
no discovery or projection implementation. Its checks include:
\begin{enumerate}
\item Both unchanged source-archive hashes and the exact selected equation
      members; every intermediate certificate is tied to its source.
\item Every Hessian-kernel identity, unit direction coordinate, unit slope,
      deleted residual, renumbering, and restoration polynomial.
\item The direct identity \eqref{eq:projected}, independently of the
      producer's shear-and-substitution algorithm.
\item The extra affine substitution by the independent identity
      $A x_p^2+B x_p+C\mapsto A S^2+B S+C$.
\item Every elementary integer row operation and its reverse replay.
\item All coefficients of each complete expanded sum of squares, and all
      dimensions, residual counts, degrees, supports, and heights.
\item A separate parser for the printed principal residuals, so the
      appendix is checked against the literal coefficient data.
\end{enumerate}
The projection checker rejects eight kinds of deliberate corruptions;
the affine checker rejects five, and the row checker rejects four.
Representative checks also passed with Python optimization enabled, so
correctness gates do not depend on removable \texttt{assert} statements.
The full package verification is reproducible with
\texttt{python3 -B verify.py}.

These finite checks establish the supplied algebraic certificates. The
cofinality theorem and the inherited arithmetic/logical calibration are
unbounded mathematical arguments. The release does not call samples an
unbounded proof or call a Python program a proof-assistant kernel.

\subsection{Two useful but limited rank facts}

For the principal system, the stacked Hessian has full column rank 234
modulo the explicitly verified prime 1,000,003. A nonzero minor modulo a
prime is nonzero over $\Q$, so its common Hessian radical over $\Q$ is
zero. Thus no further nonzero constant direction makes \emph{every}
residual have a coordinate-independent slope.

Separately, the homogeneous quadratic parts of its 171 residuals have
row rank 171 over $\Q$. Hence no nonzero constant linear combination of
the residuals is affine. The 235-variable companion has the corresponding
full row rank 170. These facts certify exhaustion of the two particular
linear-algebra moves on the displayed presentations.

They do not establish a minimum under arbitrary transformations. For a
small counterexample, take
\[
 R_1=x^2-1,\qquad R_2=xy-1,\qquad R_3=y^2-1.
\]
Their quadratic parts are independent and their common Hessian radical
is zero. Nevertheless,
\[
 yR_1-xR_2=x-y.
\]
This polynomial-multiplier identity supplies $x=y$, reducing the system
to $x^2-1=0$. Both rank checks hold before that reduction. Nonlinear
coordinate changes, polynomial residual combinations, and a different
semantic encoding therefore remain outside the limited certificate.

\section{Further research questions}

\paragraph{Polynomial residual combinations.}
Can a bounded degree-three Macaulay computation find a useful affine
consequence of the 234-variable system? The generators would be its
residuals and their products with single variables. A successful candidate
should carry an explicit polynomial-combination identity and an integral
restoration map. The counterexample above shows why this question is
strictly stronger than either rank test already completed.

\paragraph{Integral lattice generalization.}
The current implementation uses directions with a unit coordinate and a
single unit residual slope. Can Hermite or Smith normal form expose more
useful choices earlier in the source pipeline? Any nonunit invariant
factor must be retained as a divisibility condition. A rational elimination
without that accounting is insufficient for integer solutions.

\paragraph{Joint source and coordinate optimization.}
The same six common directions occur in several recent direct-source
variants. Can the source compiler use shared combinations as registers
from the outset? Doing so could avoid the later coefficient growth, and
could permit a lower-support version of the 234-variable endpoint.
This should be measured against the entire coefficient list, not a
partial gate count.

\paragraph{A smaller cofinal arithmetic family.}
The family $F^2+u^2+v^2$ handles every odd prime power and uses the already
available two-adic valuation of $F$. Is there a one-witness family with the
same uniform divisibility and cofinality properties? The tempting family
$F^2+u^2$ cannot generally replace it: if a prime $p\equiv3\pmod4$ does
not divide $F$, the congruence $u^2\equiv-F^2\pmod p$ is impossible.
Any proposed one-witness replacement must explain how it avoids this
local obstruction.

\paragraph{Formalizing the short new interface.}
The integral projection theorem is a compact first target: its proof needs
only polynomial evaluation, the constant directional-slope identity, and
integer unit arithmetic. The cofinal-spacing theorem is another contained
target, using finite-field square counts, elementary lifting, and CRT.
These are better delimited milestones than treating successful finite
coefficient checks as a completed internal ZFC development.

\section{Files and integration}

The archive includes this LaTeX source and PDF, the complete principal
residual appendix, all three final coefficient lists and expansions,
exact intermediate data, construction and independent verification code,
ordinary proof reviews, immutable source archives, and a checksum manifest.
Its intended intake location is \nolinkurl{docs/incoming}, followed by an
additive placement in the DiophantineConsistency collection. No existing
repository file was changed.

A polynomial term is encoded as \texttt{[coefficient,[variable indices]]};
repeated indices represent repeated factors and an empty index list is
constant. The polynomial is the sum of the squares of its residuals.
Variable names are explanatory metadata. The arrays and their integer
coefficients are the definition.

The new checks need only Python's standard library and no network. A user
can evaluate a supplied tuple with \nolinkurl{code/evaluate.py}. The
verification scripts do not attempt to locate a refutation or a polynomial
zero. Run \texttt{pdflatex} twice on \nolinkurl{report/report.tex} to rebuild
the PDF; typesetting versions may change PDF bytes without changing any
mathematical object.

The proof dependencies are deliberately visible. The result is a new,
fully emitted equation and a checked reduction from its documented source.
It is not a proof that this equation has no integer solution, nor a claim
that the broader search for the simplest consistency equation is finished.

\begingroup\raggedright
\begin{thebibliography}{9}
\bibitem{repo} V. Reshetnikov and ProveIt contributors,
\emph{ProveIt}, repository snapshot
\texttt{7a8da414376a80b90ecf799f37ceb26fb9ac674e}, 25 October 2026.
\url{https://github.com/VladimirReshetnikov/ProveIt/tree/7a8da414376a80b90ecf799f37ceb26fb9ac674e}.

\bibitem{overview} ProveIt,
\emph{Con(ZFC) in Diophantine form}, consolidated manuscript collection
and active proposal review, at \cite{repo}.
\url{https://github.com/VladimirReshetnikov/ProveIt/tree/7a8da414376a80b90ecf799f37ceb26fb9ac674e/Computability/DiophantineConsistency/con-zfc-diophantine-forms}.

\bibitem{shared} ProveIt incoming research,
\emph{Shared Factorial Error}, 25 October 2026,
\nolinkurl{ConZFC_Shared_Factorial_Error_2026-10-25.zip}, retained unchanged
in \nolinkurl{sources/}. Original location: \nolinkurl{docs/incoming}
at \cite{repo}. The selected parent is its
\nolinkurl{equations/equation240/equation.json}.

\bibitem{factorialfree} ProveIt incoming research,
\emph{Removing Factorial Authenticity: Binomial Prime-Power Certificates
for Diophantine Consistency}, 25 October 2026,
\nolinkurl{ConZFC_Factorial_Free_Binomial_CRT_2026-10-25.zip}, retained
unchanged in \nolinkurl{sources/}. Original location:
\nolinkurl{docs/incoming} at \cite{repo}. The selected parent is its
\nolinkurl{direct/principal/equation.json}.

\bibitem{bayer} J. Bayer, M. David, M. Ha\ss ler, Y. Matiyasevich,
and D. Schleicher,
\emph{Diophantine equations over $\Z$: universal bounds and parallel
formalization}, arXiv:2506.20909v3, 22 October 2026.
\url{https://arxiv.org/html/2506.20909v3}.

\bibitem{afp} M. David et al.,
\emph{Universal Pairs for Diophantine Equations}, Archive of Formal Proofs,
29 July 2025. \url{https://isa-afp.org/entries/Diophantine_Universal_Pairs.html}.
\end{thebibliography}
\endgroup

\clearpage\appendix
\section{Every quadratic residual of the principal polynomial}\label{app:residuals}

This appendix is the complete definition of the 234-variable equation.
All $x_0,\ldots,x_{233}$ are unrestricted integers, and
$\New=\sum_{i=0}^{170}\mathcal Q_i^2$. Coefficients and repeated variable
indices are checked against the JSON definition by a separate parser.
Line breaks in a residual have no algebraic meaning.

\begingroup\small
\setlength{\jot}{2pt}
% Generated from the literal principal equation; independently parsed by verify_printed.py.
\begin{align*}
\mathcal{Q}_{0}={}&-x_{3}^{2}+x_{4}
\end{align*}

\begin{align*}
\mathcal{Q}_{1}={}&x_{0}-x_{0}x_{2}-x_{2}+x_{9}
\end{align*}

\begin{align*}
\mathcal{Q}_{2}={}&-x_{2}x_{7}-x_{3}+x_{8}x_{9}
\end{align*}

\begin{align*}
\mathcal{Q}_{3}={}&-1+x_{1}-x_{8}x_{9}
\end{align*}

\begin{align*}
\mathcal{Q}_{4}={}&1+x_{4}-x_{5}x_{20}+x_{18}
\end{align*}

\begin{align*}
\mathcal{Q}_{5}={}&-1-x_{10}^{2}-2x_{10}x_{11}-x_{11}^{2}+x_{21}
\end{align*}

\begin{align*}
\mathcal{Q}_{6}={}&-1+x_{0}-x_{3}+x_{3}x_{10}-6x_{10}\\
&+3x_{10}^{2}-x_{10}x_{11}-x_{10}x_{12}-8x_{10}x_{21}-11x_{11}\\
&-5x_{11}^{2}+7x_{11}x_{12}-16x_{11}x_{21}+x_{12}-3x_{12}^{2}\\
&-5x_{21}+12x_{21}^{2}+x_{22}-2x_{25}+2x_{26}\\
&-x_{38}+3x_{73}-x_{78}-x_{208}^{2}-x_{209}^{2}\\
&-x_{210}^{2}-x_{211}^{2}-x_{212}^{2}+x_{218}^{2}+x_{219}^{2}
\end{align*}

\begin{align*}
\mathcal{Q}_{7}={}&-1-x_{1}+x_{8}x_{9}+2x_{10}-x_{10}^{2}\\
&-4x_{10}x_{22}-4x_{22}-4x_{22}^{2}+x_{23}
\end{align*}

\begin{align*}
\mathcal{Q}_{8}={}&-24x_{10}x_{23}+24x_{10}x_{35}+2x_{23}-4x_{23}^{2}+x_{24}\\
&-2x_{35}+4x_{35}^{2}-x_{36}
\end{align*}

\begin{align*}
\mathcal{Q}_{9}={}&-1-x_{1}+x_{8}x_{9}-x_{11}-x_{11}^{2}\\
&-2x_{11}x_{12}-x_{12}-x_{12}^{2}+x_{25}
\end{align*}

\begin{align*}
\mathcal{Q}_{10}={}&-1+x_{10}-x_{10}^{2}-2x_{10}x_{12}+x_{12}\\
&-x_{12}^{2}+x_{26}
\end{align*}

\begin{align*}
\mathcal{Q}_{11}={}&2+4x_{0}+x_{0}x_{2}+2x_{1}+x_{2}\\
&-x_{2}x_{7}-x_{3}x_{10}-x_{3}x_{11}-x_{4}+x_{5}x_{20}\\
&-x_{8}x_{9}-x_{9}+3x_{10}+7x_{10}^{2}-15x_{10}x_{11}\\
&-11x_{10}x_{12}-16x_{10}x_{21}+8x_{11}+9x_{11}^{2}-13x_{11}x_{12}\\
&-8x_{11}x_{25}-4x_{11}x_{26}-11x_{12}+4x_{12}^{2}-2x_{12}x_{13}\\
&-28x_{12}x_{21}-x_{13}^{2}-x_{18}-13x_{21}+8x_{21}^{2}\\
&-8x_{21}x_{26}+2x_{24}-11x_{25}-20x_{25}^{2}-8x_{25}x_{26}\\
&+4x_{26}-4x_{26}^{2}+x_{27}+2x_{28}+x_{29}\\
&+2x_{30}-x_{31}-4x_{31}^{2}+3x_{32}+2x_{34}\\
&+4x_{36}+2x_{37}-x_{38}+2x_{40}+x_{41}\\
&+2x_{43}+2x_{45}+x_{61}-4x_{62}^{2}+3x_{63}\\
&-7x_{73}+x_{77}+x_{78}+x_{138}x_{192}+x_{192}\\
&-x_{193}+x_{204}x_{217}-5x_{208}^{2}-x_{212}
\end{align*}

\begin{align*}
\mathcal{Q}_{12}={}&1-x_{0}+9x_{10}-4x_{10}^{2}+x_{10}x_{11}\\
&+3x_{10}x_{12}-8x_{10}x_{25}+2x_{11}^{2}+3x_{11}x_{12}-5x_{12}\\
&+4x_{12}^{2}-2x_{12}x_{13}-x_{13}^{2}+x_{21}+5x_{25}\\
&-x_{26}-4x_{27}^{2}+x_{28}-4x_{33}-4x_{33}^{2}\\
&+x_{34}+x_{41}-2x_{45}-3x_{73}+x_{208}^{2}
\end{align*}

\begin{align*}
\mathcal{Q}_{13}={}&-8x_{10}-x_{10}x_{11}+x_{10}x_{12}+8x_{10}x_{25}+8x_{10}x_{26}\\
&-6x_{11}+x_{11}^{2}+3x_{11}x_{12}+8x_{11}x_{25}+8x_{11}x_{26}\\
&-4x_{21}+7x_{25}-4x_{25}^{2}-8x_{25}x_{26}+8x_{26}\\
&-4x_{26}^{2}+x_{29}+x_{73}
\end{align*}

\begin{align*}
\mathcal{Q}_{14}={}&-1-x_{1}+x_{3}x_{11}+x_{8}x_{9}-x_{10}\\
&+6x_{10}^{2}+7x_{10}x_{11}+x_{10}x_{12}+8x_{10}x_{21}-8x_{10}x_{26}\\
&-8x_{10}x_{29}+4x_{11}-14x_{11}^{2}+11x_{11}x_{12}-16x_{11}x_{21}\\
&+8x_{11}x_{25}-24x_{11}x_{26}-16x_{11}x_{29}-4x_{11}x_{49}-x_{12}\\
&+6x_{12}^{2}-2x_{12}x_{13}+8x_{12}x_{26}+4x_{12}x_{29}-8x_{12}x_{49}\\
&-x_{13}^{2}+3x_{21}-4x_{21}^{2}-16x_{21}x_{26}-8x_{21}x_{29}\\
&+8x_{25}-4x_{25}^{2}-8x_{25}x_{26}+15x_{26}-4x_{26}^{2}\\
&-4x_{27}^{2}+x_{28}+5x_{29}+x_{30}+x_{41}\\
&-x_{44}-2x_{45}-4x_{49}^{2}+x_{50}-x_{73}\\
&-x_{77}
\end{align*}

\begin{align*}
\mathcal{Q}_{15}={}&-1+2x_{0}-x_{1}+x_{8}x_{9}+7x_{10}\\
&+9x_{10}^{2}-12x_{10}x_{11}-10x_{10}x_{12}-12x_{10}x_{25}+16x_{10}x_{26}\\
&+24x_{10}x_{39}+4x_{10}x_{54}-12x_{11}+16x_{11}^{2}+2x_{11}x_{12}\\
&-8x_{11}x_{25}+4x_{11}x_{26}-20x_{12}-2x_{12}^{2}-16x_{12}x_{25}\\
&-8x_{12}x_{26}+16x_{12}x_{39}+8x_{12}x_{54}+2x_{21}-19x_{25}\\
&-12x_{25}^{2}-8x_{25}x_{26}+3x_{26}-4x_{26}^{2}+8x_{26}x_{39}\\
&-8x_{27}^{2}+2x_{28}+x_{29}+x_{31}+8x_{33}\\
&+8x_{33}^{2}-2x_{34}-4x_{39}+4x_{39}^{2}-x_{40}\\
&-3x_{45}+4x_{54}+4x_{54}^{2}-x_{55}-16x_{73}\\
&-2x_{208}^{2}
\end{align*}

\begin{align*}
\mathcal{Q}_{16}={}&-2-x_{0}-2x_{1}+x_{2}x_{7}+x_{3}x_{10}\\
&+x_{3}x_{11}+x_{4}-x_{5}x_{20}+x_{8}x_{9}+9x_{10}\\
&+14x_{10}^{2}-x_{10}x_{11}+5x_{10}x_{12}-8x_{10}x_{25}+16x_{10}x_{26}\\
&+24x_{10}x_{39}+5x_{11}-x_{11}^{2}-x_{11}x_{12}-8x_{11}x_{25}\\
&-8x_{11}x_{26}+6x_{12}+8x_{12}x_{25}+16x_{12}x_{26}+16x_{12}x_{39}\\
&+x_{18}+7x_{21}-2x_{25}+8x_{25}x_{26}+5x_{26}\\
&+8x_{26}^{2}+8x_{26}x_{39}-x_{29}-8x_{33}-8x_{33}^{2}\\
&+2x_{34}-4x_{39}+4x_{39}^{2}-x_{40}+x_{45}\\
&-5x_{73}-x_{77}-x_{78}+x_{208}^{2}
\end{align*}

\begin{align*}
\mathcal{Q}_{17}={}&-1+2x_{0}-x_{1}+x_{3}x_{11}+x_{8}x_{9}\\
&+x_{10}+5x_{10}^{2}-2x_{10}x_{11}+6x_{10}x_{12}+4x_{10}x_{25}\\
&+16x_{10}x_{26}-5x_{11}+x_{11}^{2}-2x_{11}x_{12}-16x_{11}x_{25}\\
&+4x_{11}x_{26}-2x_{12}-2x_{12}^{2}+24x_{12}x_{26}+2x_{21}\\
&-14x_{25}-16x_{25}^{2}+16x_{26}+4x_{26}^{2}+4x_{27}^{2}\\
&-x_{28}-x_{31}-4x_{31}^{2}+x_{32}-4x_{33}\\
&-4x_{33}^{2}+x_{34}+5x_{45}-2x_{73}-x_{77}\\
&-2x_{208}^{2}-x_{211}^{2}-x_{212}^{2}+x_{219}^{2}
\end{align*}

\begin{align*}
\mathcal{Q}_{18}={}&1-2x_{10}+4x_{10}^{2}+4x_{10}x_{11}-4x_{10}x_{12}\\
&-8x_{10}x_{26}-8x_{10}x_{61}-20x_{11}-7x_{11}^{2}+14x_{11}x_{12}\\
&+8x_{11}x_{21}+16x_{11}x_{25}-8x_{11}x_{26}-16x_{11}x_{61}-5x_{12}^{2}\\
&+2x_{12}x_{13}+4x_{12}x_{59}+x_{13}^{2}+7x_{21}-8x_{21}x_{25}\\
&+8x_{21}x_{61}+2x_{25}+8x_{25}x_{26}-4x_{26}+4x_{26}^{2}\\
&-x_{29}+x_{33}-x_{41}+4x_{59}^{2}-x_{60}\\
&-4x_{61}+4x_{61}^{2}+6x_{73}
\end{align*}

\begin{align*}
\mathcal{Q}_{19}={}&-1-x_{1}+x_{8}x_{9}-5x_{10}-x_{10}^{2}\\
&-2x_{10}x_{12}+x_{11}+x_{11}^{2}+2x_{11}x_{12}-2x_{12}\\
&-x_{25}-x_{26}-4x_{33}-4x_{33}^{2}+x_{34}
\end{align*}

\begin{align*}
\mathcal{Q}_{20}={}&-1+x_{0}+x_{3}x_{11}+18x_{10}+12x_{10}^{2}\\
&+11x_{10}x_{11}-5x_{10}x_{12}-8x_{10}x_{21}+16x_{10}x_{22}+8x_{10}x_{25}\\
&+8x_{10}x_{26}+8x_{10}x_{29}-4x_{10}x_{41}+5x_{11}-12x_{11}^{2}\\
&+5x_{11}x_{12}-8x_{11}x_{21}+64x_{11}x_{22}+16x_{11}x_{25}+24x_{11}x_{26}\\
&+16x_{11}x_{29}-14x_{12}+6x_{12}^{2}+4x_{12}x_{21}+8x_{12}x_{26}\\
&-4x_{12}x_{29}-8x_{12}x_{41}+17x_{21}-20x_{21}^{2}+32x_{21}x_{22}\\
&-8x_{21}x_{25}+16x_{21}x_{26}+8x_{21}x_{29}+2x_{22}-16x_{22}^{2}\\
&+3x_{23}-2x_{24}-3x_{25}+4x_{25}^{2}+x_{26}\\
&+4x_{27}^{2}-3x_{28}-5x_{29}-3x_{30}-2x_{32}\\
&-3x_{33}-4x_{33}^{2}-x_{34}+x_{35}-4x_{36}\\
&-2x_{37}+x_{38}-2x_{40}+4x_{41}-4x_{41}^{2}\\
&+x_{42}-2x_{43}+x_{44}-x_{61}-4x_{62}\\
&-2x_{63}-5x_{73}-x_{77}-x_{204}x_{217}-x_{208}^{2}\\
&+x_{212}
\end{align*}

\begin{align*}
\mathcal{Q}_{21}={}&2+x_{0}+x_{1}-x_{2}x_{7}-x_{3}x_{10}\\
&-x_{3}x_{11}-x_{4}+x_{5}x_{20}-7x_{10}-24x_{10}^{2}\\
&+7x_{10}x_{11}+5x_{10}x_{12}+8x_{10}x_{25}-24x_{10}x_{35}-x_{11}\\
&+x_{11}^{2}-5x_{11}x_{12}-2x_{12}^{2}-x_{18}-5x_{21}\\
&-4x_{25}-3x_{26}+4x_{27}^{2}-x_{28}+4x_{33}\\
&+4x_{33}^{2}-x_{34}+2x_{35}-4x_{35}^{2}+x_{36}\\
&+2x_{45}+3x_{73}+x_{77}+x_{78}-x_{208}^{2}
\end{align*}

\begin{align*}
\mathcal{Q}_{22}={}&-x_{11}^{2}-4x_{11}x_{36}-4x_{36}^{2}+x_{37}
\end{align*}

\begin{align*}
\mathcal{Q}_{23}={}&-1-x_{1}+x_{2}x_{7}+x_{3}x_{10}+x_{3}x_{11}\\
&+x_{4}-x_{5}x_{20}-x_{10}-x_{10}^{2}-2x_{10}x_{11}\\
&-2x_{11}+2x_{11}x_{12}+x_{12}+x_{12}^{2}+x_{18}\\
&-9x_{21}-16x_{21}^{2}+2x_{24}-x_{25}+2x_{28}\\
&+2x_{30}+2x_{32}+2x_{34}+4x_{36}+2x_{37}\\
&+x_{38}+2x_{40}+2x_{43}+4x_{62}+2x_{63}\\
&-x_{77}-x_{78}+x_{204}x_{217}-x_{212}
\end{align*}

\begin{align*}
\mathcal{Q}_{24}={}&1+x_{0}+x_{10}-x_{10}^{2}+5x_{10}x_{11}\\
&+5x_{10}x_{12}+8x_{10}x_{26}+18x_{11}-13x_{11}^{2}-7x_{11}x_{12}\\
&-16x_{11}x_{25}+8x_{11}x_{26}-24x_{11}x_{38}+7x_{12}+x_{12}^{2}\\
&-8x_{12}x_{25}-16x_{12}x_{38}-6x_{21}+5x_{25}-12x_{25}^{2}\\
&-8x_{25}x_{26}-8x_{25}x_{38}+7x_{26}-4x_{26}^{2}-4x_{27}^{2}\\
&+x_{28}+x_{29}-4x_{33}-4x_{33}^{2}+x_{34}\\
&+4x_{38}-4x_{38}^{2}+x_{39}-x_{45}-x_{73}\\
&-x_{208}^{2}
\end{align*}

\begin{align*}
\mathcal{Q}_{25}={}&-x_{12}^{2}-2x_{12}x_{13}-x_{13}^{2}+x_{41}
\end{align*}

\begin{align*}
\mathcal{Q}_{26}={}&-1-x_{1}+x_{8}x_{9}-4x_{10}+5x_{10}^{2}\\
&+2x_{10}x_{11}-6x_{10}x_{12}-8x_{10}x_{26}-4x_{10}x_{41}+5x_{11}\\
&+x_{11}^{2}-8x_{11}x_{25}-8x_{11}x_{26}-13x_{12}+3x_{12}^{2}\\
&+16x_{12}x_{26}-8x_{12}x_{41}+6x_{21}-11x_{25}+8x_{25}^{2}\\
&+12x_{26}+4x_{27}^{2}-x_{28}-x_{29}+4x_{33}\\
&+4x_{33}^{2}-x_{34}+4x_{41}-4x_{41}^{2}+x_{42}\\
&-6x_{73}
\end{align*}

\begin{align*}
\mathcal{Q}_{27}={}&1+x_{1}-x_{3}x_{11}-x_{8}x_{9}+x_{10}^{2}\\
&-x_{10}x_{11}-x_{10}x_{12}-8x_{10}x_{41}-x_{11}-x_{11}^{2}\\
&-x_{11}x_{12}+x_{12}-2x_{12}^{2}+2x_{12}x_{13}-16x_{12}x_{41}\\
&+x_{13}^{2}+x_{21}-4x_{26}-16x_{26}^{2}-32x_{26}x_{41}\\
&-16x_{26}x_{42}+3x_{41}-24x_{41}^{2}-16x_{41}x_{42}-4x_{42}^{2}\\
&+x_{43}-x_{73}+x_{77}
\end{align*}

\begin{align*}
\mathcal{Q}_{28}={}&1+x_{0}+x_{1}-x_{8}x_{9}-x_{10}\\
&-10x_{10}^{2}+6x_{10}x_{11}+2x_{10}x_{12}+24x_{10}x_{26}+16x_{10}x_{29}\\
&+5x_{11}-6x_{11}x_{12}-8x_{11}x_{25}-8x_{11}x_{26}+x_{12}\\
&-6x_{12}^{2}+2x_{12}x_{13}+8x_{12}x_{25}-8x_{12}x_{26}-4x_{12}x_{29}\\
&+x_{13}^{2}+x_{21}-5x_{25}+8x_{25}x_{26}-7x_{26}\\
&-12x_{26}^{2}-16x_{26}x_{29}+4x_{27}^{2}-x_{28}-x_{29}\\
&-4x_{29}^{2}+4x_{33}+4x_{33}^{2}-x_{34}-x_{41}\\
&+x_{44}+3x_{45}-x_{208}^{2}
\end{align*}

\begin{align*}
\mathcal{Q}_{29}={}&1+x_{1}-x_{8}x_{9}-x_{10}^{2}-2x_{10}x_{12}\\
&+x_{11}^{2}+2x_{11}x_{12}-4x_{12}+8x_{12}x_{25}+x_{21}\\
&+6x_{25}-4x_{25}^{2}+x_{45}-2x_{73}
\end{align*}

\begin{align*}
\mathcal{Q}_{30}={}&-1-x_{2}-x_{6}+x_{6}^{2}-x_{10}\\
&+2x_{10}^{2}-x_{10}x_{11}+5x_{10}x_{12}-x_{11}-3x_{11}^{2}\\
&-5x_{11}x_{12}+8x_{12}-6x_{12}^{2}+2x_{12}x_{13}+8x_{12}x_{25}\\
&+12x_{12}x_{45}+x_{13}^{2}-x_{21}-12x_{25}-8x_{25}^{2}\\
&-16x_{25}x_{45}-x_{26}-x_{41}-2x_{45}-4x_{45}^{2}\\
&+x_{46}+3x_{73}-x_{138}x_{192}-x_{192}+x_{193}\\
&-x_{209}^{2}-x_{210}^{2}-x_{211}^{2}-x_{212}^{2}+x_{218}^{2}\\
&+x_{219}^{2}
\end{align*}

\begin{align*}
\mathcal{Q}_{31}={}&1+x_{1}-x_{3}x_{11}-x_{8}x_{9}+2x_{11}\\
&-2x_{11}x_{12}-4x_{11}x_{46}+x_{12}-3x_{12}^{2}-8x_{12}x_{46}\\
&-x_{25}-4x_{46}^{2}+x_{47}+x_{77}
\end{align*}

\begin{align*}
\mathcal{Q}_{32}={}&-1-4x_{47}-4x_{47}^{2}+x_{48}
\end{align*}

\begin{align*}
\mathcal{Q}_{33}={}&1-x_{12}^{2}-8x_{12}x_{47}-4x_{12}x_{48}+8x_{47}\\
&-8x_{47}^{2}-16x_{47}x_{48}-2x_{48}-4x_{48}^{2}+x_{49}
\end{align*}

\begin{align*}
\mathcal{Q}_{34}={}&2+x_{1}-x_{3}x_{11}-x_{8}x_{9}+2x_{11}\\
&-2x_{11}x_{12}-4x_{11}x_{49}+x_{12}-3x_{12}^{2}-8x_{12}x_{49}\\
&-x_{25}-4x_{49}^{2}+x_{50}+x_{77}
\end{align*}

\begin{align*}
\mathcal{Q}_{35}={}&-4x_{44}^{2}-8x_{44}x_{50}-4x_{50}^{2}+x_{51}
\end{align*}

\begin{align*}
\mathcal{Q}_{36}={}&-2-x_{1}+x_{2}x_{7}+x_{3}x_{10}+x_{3}x_{11}\\
&+x_{4}-x_{5}x_{20}+6x_{10}+2x_{10}x_{12}-8x_{10}x_{44}\\
&-4x_{10}x_{51}+x_{11}+x_{11}^{2}+2x_{11}x_{12}+2x_{12}^{2}\\
&+x_{18}-x_{25}-x_{26}+24x_{44}-12x_{44}^{2}\\
&+8x_{44}x_{50}-16x_{44}x_{51}+4x_{50}^{2}+11x_{51}-4x_{51}^{2}\\
&+x_{52}-x_{77}-x_{78}
\end{align*}

\begin{align*}
\mathcal{Q}_{37}={}&-1-4x_{52}-4x_{52}^{2}+x_{53}
\end{align*}

\begin{align*}
\mathcal{Q}_{38}={}&1+x_{1}-x_{8}x_{9}-x_{12}^{2}-8x_{12}x_{52}\\
&-4x_{12}x_{53}+8x_{52}-8x_{52}^{2}-16x_{52}x_{53}-2x_{53}\\
&-4x_{53}^{2}+x_{54}
\end{align*}

\begin{align*}
\mathcal{Q}_{39}={}&1+x_{1}-x_{8}x_{9}-2x_{10}-x_{10}^{2}\\
&-x_{10}x_{11}-3x_{10}x_{12}-4x_{10}x_{54}+x_{11}x_{12}-3x_{12}\\
&-x_{12}^{2}-8x_{12}x_{54}+x_{21}-x_{25}-x_{26}\\
&-4x_{54}-4x_{54}^{2}+x_{55}-x_{73}
\end{align*}

\begin{align*}
\mathcal{Q}_{40}={}&2-x_{0}+x_{1}-x_{2}x_{7}-x_{3}x_{10}\\
&-x_{3}x_{11}-x_{4}+x_{5}x_{20}-x_{18}-12x_{55}\\
&-4x_{55}^{2}+x_{56}+x_{77}+x_{78}+x_{208}^{2}
\end{align*}

\begin{align*}
\mathcal{Q}_{41}={}&-x_{10}x_{11}+x_{10}x_{12}-x_{11}^{2}-x_{11}x_{12}+2x_{12}\\
&-2x_{12}^{2}+16x_{12}x_{55}+8x_{12}x_{56}-3x_{25}-16x_{25}x_{55}\\
&-8x_{25}x_{56}-x_{45}+8x_{55}-16x_{55}^{2}-16x_{55}x_{56}\\
&+4x_{56}-4x_{56}^{2}+x_{57}+x_{73}
\end{align*}

\begin{align*}
\mathcal{Q}_{42}={}&1-x_{0}+x_{3}-x_{3}x_{10}-x_{3}x_{11}\\
&-x_{10}+x_{10}^{2}+2x_{10}x_{12}-2x_{11}-4x_{11}^{2}\\
&+2x_{11}x_{12}-8x_{11}x_{25}-4x_{11}x_{57}+4x_{12}-4x_{12}^{2}\\
&+2x_{12}x_{13}-12x_{12}x_{45}+16x_{12}x_{57}+x_{13}^{2}-5x_{25}\\
&+4x_{25}^{2}+16x_{25}x_{45}-16x_{25}x_{57}-x_{26}-x_{41}\\
&-x_{45}+4x_{45}^{2}-x_{46}-12x_{50}-4x_{50}^{2}\\
&-4x_{57}^{2}+x_{58}+x_{77}+x_{78}+x_{208}^{2}
\end{align*}

\begin{align*}
\mathcal{Q}_{43}={}&-2+x_{0}+x_{0}x_{2}-2x_{1}+x_{2}x_{7}\\
&+x_{3}^{2}+x_{3}x_{10}+x_{3}x_{11}-x_{5}x_{20}-x_{6}\\
&+x_{6}^{2}+x_{8}x_{9}-x_{9}+x_{18}+24x_{50}\\
&-8x_{50}^{2}-16x_{50}x_{58}-2x_{58}-4x_{58}^{2}+x_{59}\\
&-x_{77}-x_{78}-x_{138}x_{192}-x_{192}+x_{193}\\
&-2x_{208}^{2}-x_{209}^{2}-x_{210}^{2}-x_{211}^{2}-x_{212}^{2}\\
&+x_{218}^{2}+x_{219}^{2}
\end{align*}

\begin{align*}
\mathcal{Q}_{44}={}&2+x_{1}-x_{2}x_{7}-x_{3}x_{10}-x_{3}x_{11}\\
&-x_{4}+x_{5}x_{20}-7x_{10}+x_{10}^{2}-x_{10}x_{11}\\
&-x_{10}x_{12}+8x_{10}x_{61}-9x_{11}+4x_{11}^{2}+x_{11}x_{12}\\
&+16x_{11}x_{61}-3x_{12}^{2}+2x_{12}x_{13}+4x_{12}x_{21}-4x_{12}x_{59}\\
&+x_{13}^{2}-x_{18}+6x_{21}-8x_{21}x_{61}-2x_{25}\\
&+x_{26}-x_{41}-4x_{59}^{2}+x_{60}+3x_{61}\\
&-4x_{61}^{2}+5x_{73}+x_{77}+x_{78}
\end{align*}

\begin{align*}
\mathcal{Q}_{45}={}&2+x_{1}-x_{3}x_{11}-x_{8}x_{9}+4x_{10}\\
&+2x_{10}^{2}+8x_{10}x_{21}+3x_{11}-11x_{11}^{2}+2x_{11}x_{12}\\
&+16x_{11}x_{21}+2x_{12}+2x_{12}^{2}-4x_{12}x_{21}-10x_{21}\\
&-4x_{21}^{2}+x_{25}+x_{61}-4x_{73}+x_{77}
\end{align*}

\begin{align*}
\mathcal{Q}_{46}={}&1+x_{1}-x_{8}x_{9}+2x_{10}+x_{10}^{2}\\
&-3x_{10}x_{11}-3x_{10}x_{12}+8x_{10}x_{21}-8x_{10}x_{25}+16x_{10}x_{60}\\
&-3x_{11}+x_{11}^{2}-3x_{11}x_{12}+16x_{11}x_{21}+32x_{11}x_{60}\\
&-6x_{12}+5x_{12}^{2}-2x_{12}x_{13}+12x_{12}x_{21}+8x_{12}x_{59}\\
&-x_{13}^{2}-x_{21}-4x_{21}^{2}-16x_{21}x_{60}-4x_{27}^{2}\\
&+x_{28}-4x_{33}-4x_{33}^{2}+x_{34}+x_{41}\\
&-2x_{45}+8x_{59}^{2}-2x_{60}-4x_{60}^{2}-3x_{61}\\
&+x_{62}+5x_{73}
\end{align*}

\begin{align*}
\mathcal{Q}_{47}={}&-1-x_{1}+x_{8}x_{9}-4x_{62}-4x_{62}^{2}\\
&+x_{63}
\end{align*}

\begin{align*}
\mathcal{Q}_{48}={}&-x_{3}x_{15}+x_{64}
\end{align*}

\begin{align*}
\mathcal{Q}_{49}={}&x_{14}-x_{15}-2x_{24}x_{64}-x_{65}
\end{align*}

\begin{align*}
\mathcal{Q}_{50}={}&x_{14}-x_{15}-2x_{28}x_{64}-x_{66}
\end{align*}

\begin{align*}
\mathcal{Q}_{51}={}&-x_{3}x_{15}-2x_{30}x_{64}+2x_{43}x_{64}-x_{64}-x_{67}\\
&+x_{72}
\end{align*}

\begin{align*}
\mathcal{Q}_{52}={}&x_{14}-x_{15}-2x_{32}x_{64}-x_{68}
\end{align*}

\begin{align*}
\mathcal{Q}_{53}={}&x_{14}-x_{15}+2x_{24}x_{64}-2x_{30}x_{64}-2x_{34}x_{64}\\
&+x_{65}-x_{67}-x_{69}
\end{align*}

\begin{align*}
\mathcal{Q}_{54}={}&x_{14}-x_{15}-4x_{36}x_{64}-2x_{37}x_{64}-x_{70}
\end{align*}

\begin{align*}
\mathcal{Q}_{55}={}&x_{14}-x_{15}-2x_{40}x_{64}-x_{71}
\end{align*}

\begin{align*}
\mathcal{Q}_{56}={}&-2x_{3}x_{15}+x_{14}-x_{15}-2x_{43}x_{64}-2x_{64}\\
&-x_{72}
\end{align*}

\begin{align*}
\mathcal{Q}_{57}={}&-x_{10}x_{11}+x_{10}x_{12}+x_{11}+x_{11}x_{12}-x_{12}\\
&-x_{12}^{2}+x_{73}
\end{align*}

\begin{align*}
\mathcal{Q}_{58}={}&1+2x_{13}+x_{13}^{2}-x_{73}^{2}+x_{74}
\end{align*}

\begin{align*}
\mathcal{Q}_{59}={}&x_{3}x_{15}-4x_{62}x_{64}-2x_{63}x_{64}+x_{64}+x_{75}
\end{align*}

\begin{align*}
\mathcal{Q}_{60}={}&-x_{14}^{2}+2x_{14}x_{15}+2x_{14}x_{75}-x_{15}^{2}-2x_{15}x_{75}\\
&-x_{74}^{2}-x_{75}^{2}+x_{76}
\end{align*}

\begin{align*}
\mathcal{Q}_{61}={}&-1-x_{3}x_{11}+x_{77}
\end{align*}

\begin{align*}
\mathcal{Q}_{62}={}&-x_{15}x_{77}+x_{79}x_{80}+x_{81}
\end{align*}

\begin{align*}
\mathcal{Q}_{63}={}&-1+x_{3}-x_{3}x_{10}+x_{78}
\end{align*}

\begin{align*}
\mathcal{Q}_{64}={}&-2-x_{2}x_{7}+x_{3}+2x_{3}x_{10}+3x_{3}x_{11}\\
&-2x_{3}x_{21}+x_{8}x_{9}+x_{77}+x_{79}
\end{align*}

\begin{align*}
\mathcal{Q}_{65}={}&-x_{14}x_{78}+x_{80}
\end{align*}

\begin{align*}
\mathcal{Q}_{66}={}&5x_{14}+x_{14}x_{77}+7x_{15}-5x_{15}x_{77}-4x_{15}x_{78}\\
&-x_{15}x_{79}-2x_{24}x_{64}-2x_{28}x_{64}-2x_{32}x_{64}-2x_{34}x_{64}\\
&-2x_{40}x_{64}-x_{65}-x_{66}-x_{68}-x_{69}\\
&-x_{71}-3x_{79}x_{80}-3x_{81}+x_{82}
\end{align*}

\begin{align*}
\mathcal{Q}_{67}={}&-x_{14}x_{65}+x_{15}x_{65}+x_{83}
\end{align*}

\begin{align*}
\mathcal{Q}_{68}={}&-x_{66}x_{83}+x_{84}
\end{align*}

\begin{align*}
\mathcal{Q}_{69}={}&-x_{67}x_{84}+x_{85}
\end{align*}

\begin{align*}
\mathcal{Q}_{70}={}&-x_{68}x_{85}+x_{86}
\end{align*}

\begin{align*}
\mathcal{Q}_{71}={}&-x_{69}x_{86}+x_{87}
\end{align*}

\begin{align*}
\mathcal{Q}_{72}={}&-x_{70}x_{87}+x_{88}
\end{align*}

\begin{align*}
\mathcal{Q}_{73}={}&-x_{71}x_{88}+x_{89}
\end{align*}

\begin{align*}
\mathcal{Q}_{74}={}&-x_{72}x_{89}+x_{90}
\end{align*}

\begin{align*}
\mathcal{Q}_{75}={}&-x_{81}x_{90}+x_{91}
\end{align*}

\begin{align*}
\mathcal{Q}_{76}={}&-x_{82}x_{91}+x_{92}
\end{align*}

\begin{align*}
\mathcal{Q}_{77}={}&-x_{92}^{2}+x_{93}
\end{align*}

\begin{align*}
\mathcal{Q}_{78}={}&-x_{2}-x_{2}x_{18}+x_{18}+x_{94}
\end{align*}

\begin{align*}
\mathcal{Q}_{79}={}&-x_{16}x_{94}+x_{95}
\end{align*}

\begin{align*}
\mathcal{Q}_{80}={}&-2x_{1}^{2}+2x_{1}x_{14}+2x_{1}x_{15}+2x_{1}x_{95}+2x_{1}x_{96}\\
&-x_{14}^{2}-2x_{14}x_{95}-x_{15}^{2}-2x_{15}x_{96}+x_{19}x_{97}\\
&-x_{76}x_{93}-x_{95}^{2}-x_{96}^{2}
\end{align*}

\begin{align*}
\mathcal{Q}_{81}={}&-x_{2}x_{17}+x_{17}-x_{17}x_{94}+x_{96}
\end{align*}

\begin{align*}
\mathcal{Q}_{82}={}&-x_{100}-2x_{100}x_{231}+x_{102}
\end{align*}

\begin{align*}
\mathcal{Q}_{83}={}&-x_{5}x_{124}+x_{98}x_{222}-x_{101}-x_{102}x_{124}
\end{align*}

\begin{align*}
\mathcal{Q}_{84}={}&x_{103}x_{223}-x_{105}x_{127}-x_{106}
\end{align*}

\begin{align*}
\mathcal{Q}_{85}={}&x_{107}x_{223}-x_{108}x_{127}-x_{109}
\end{align*}

\begin{align*}
\mathcal{Q}_{86}={}&x_{110}x_{223}-x_{111}x_{127}-x_{112}
\end{align*}

\begin{align*}
\mathcal{Q}_{87}={}&x_{113}x_{223}-x_{114}x_{127}-x_{115}
\end{align*}

\begin{align*}
\mathcal{Q}_{88}={}&x_{116}x_{224}-x_{118}x_{127}-x_{119}
\end{align*}

\begin{align*}
\mathcal{Q}_{89}={}&x_{120}x_{224}-x_{121}x_{127}-x_{122}
\end{align*}

\begin{align*}
\mathcal{Q}_{90}={}&-x_{19}x_{124}+x_{127}
\end{align*}

\begin{align*}
\mathcal{Q}_{91}={}&x_{123}-x_{125}x_{127}-x_{126}
\end{align*}

\begin{align*}
\mathcal{Q}_{92}={}&-x_{127}x_{129}+x_{128}-x_{130}
\end{align*}

\begin{align*}
\mathcal{Q}_{93}={}&-x_{127}x_{132}+x_{131}x_{222}-x_{133}
\end{align*}

\begin{align*}
\mathcal{Q}_{94}={}&-x_{127}x_{135}+x_{134}x_{223}-x_{136}
\end{align*}

\begin{align*}
\mathcal{Q}_{95}={}&x_{137}-2x_{137}x_{139}+x_{140}
\end{align*}

\begin{align*}
\mathcal{Q}_{96}={}&-x_{137}x_{140}-x_{140}+x_{141}
\end{align*}

\begin{align*}
\mathcal{Q}_{97}={}&x_{140}^{2}-x_{141}^{2}+x_{161}
\end{align*}

\begin{align*}
\mathcal{Q}_{98}={}&-x_{138}x_{143}-x_{143}+x_{144}
\end{align*}

\begin{align*}
\mathcal{Q}_{99}={}&2x_{142}+x_{142}^{2}+x_{143}^{2}-x_{144}^{2}
\end{align*}

\begin{align*}
\mathcal{Q}_{100}={}&x_{4}^{2}-2x_{4}x_{138}+x_{138}^{2}+x_{145}-x_{161}
\end{align*}

\begin{align*}
\mathcal{Q}_{101}={}&-1-x_{4}x_{143}+x_{6}+x_{138}x_{143}-x_{142}\\
&+x_{145}x_{146}
\end{align*}

\begin{align*}
\mathcal{Q}_{102}={}&1-2x_{5}+x_{5}^{2}-2x_{5}x_{138}+x_{147}
\end{align*}

\begin{align*}
\mathcal{Q}_{103}={}&-x_{5}x_{192}+x_{19}+x_{147}x_{148}-x_{191}+x_{193}
\end{align*}

\begin{align*}
\mathcal{Q}_{104}={}&-1+2x_{138}x_{207}-x_{149}+2x_{207}-x_{207}^{2}
\end{align*}

\begin{align*}
\mathcal{Q}_{105}={}&1+x_{149}x_{150}-x_{191}-x_{192}x_{207}+x_{193}\\
&+2x_{231}
\end{align*}

\begin{align*}
\mathcal{Q}_{106}={}&-x_{138}x_{152}-x_{152}+x_{153}
\end{align*}

\begin{align*}
\mathcal{Q}_{107}={}&2x_{151}+x_{151}^{2}+x_{152}^{2}-x_{153}^{2}
\end{align*}

\begin{align*}
\mathcal{Q}_{108}={}&-x_{138}+x_{138}^{2}-4x_{138}x_{231}-x_{149}x_{150}+x_{154}\\
&-x_{161}+x_{191}+x_{192}x_{207}-x_{193}+x_{209}x_{215}\\
&-x_{213}+2x_{231}+4x_{231}^{2}
\end{align*}

\begin{align*}
\mathcal{Q}_{109}={}&-1+x_{98}+x_{138}x_{152}-x_{151}-x_{152}\\
&-2x_{152}x_{231}+x_{154}x_{155}
\end{align*}

\begin{align*}
\mathcal{Q}_{110}={}&-1+x_{99}+x_{138}x_{143}-2x_{138}x_{156}-x_{142}\\
&-2x_{143}x_{231}+x_{154}x_{156}+x_{156}+4x_{156}x_{231}
\end{align*}

\begin{align*}
\mathcal{Q}_{111}={}&1+x_{138}-x_{138}x_{158}-x_{158}+x_{159}
\end{align*}

\begin{align*}
\mathcal{Q}_{112}={}&-x_{138}+x_{138}x_{158}+2x_{157}+x_{157}^{2}-x_{158}\\
&+x_{158}^{2}-x_{159}-x_{159}^{2}
\end{align*}

\begin{align*}
\mathcal{Q}_{113}={}&x_{103}-x_{138}+x_{138}x_{158}+x_{154}x_{160}-x_{157}\\
&-x_{158}-2x_{158}x_{231}+2x_{231}
\end{align*}

\begin{align*}
\mathcal{Q}_{114}={}&-x_{138}-x_{138}^{2}+x_{161}+x_{209}x_{215}-x_{213}
\end{align*}

\begin{align*}
\mathcal{Q}_{115}={}&x_{104}+x_{138}x_{144}-2x_{138}x_{162}-x_{142}-2x_{142}x_{231}\\
&-x_{143}x_{161}-2x_{144}x_{231}+x_{149}x_{150}+x_{154}x_{162}+x_{162}\\
&+4x_{162}x_{231}-x_{191}-x_{192}x_{207}+x_{193}
\end{align*}

\begin{align*}
\mathcal{Q}_{116}={}&-x_{138}x_{164}-x_{164}+x_{165}
\end{align*}

\begin{align*}
\mathcal{Q}_{117}={}&2x_{163}+x_{163}^{2}+x_{164}^{2}-x_{165}^{2}
\end{align*}

\begin{align*}
\mathcal{Q}_{118}={}&-1+x_{107}+x_{138}x_{164}+x_{154}x_{166}-x_{163}\\
&-x_{164}-2x_{164}x_{231}
\end{align*}

\begin{align*}
\mathcal{Q}_{119}={}&-x_{138}x_{168}-x_{168}+x_{169}
\end{align*}

\begin{align*}
\mathcal{Q}_{120}={}&2x_{167}+x_{167}^{2}+x_{168}^{2}-x_{169}^{2}
\end{align*}

\begin{align*}
\mathcal{Q}_{121}={}&-1+x_{110}+x_{138}x_{168}+x_{154}x_{170}-x_{167}\\
&-x_{168}-2x_{168}x_{231}
\end{align*}

\begin{align*}
\mathcal{Q}_{122}={}&-x_{138}x_{172}-x_{172}+x_{173}
\end{align*}

\begin{align*}
\mathcal{Q}_{123}={}&2x_{171}+x_{171}^{2}+x_{172}^{2}-x_{173}^{2}
\end{align*}

\begin{align*}
\mathcal{Q}_{124}={}&-1+x_{113}+x_{138}x_{172}+x_{154}x_{174}-x_{171}\\
&-x_{172}-2x_{172}x_{231}
\end{align*}

\begin{align*}
\mathcal{Q}_{125}={}&-x_{138}x_{176}-x_{176}+x_{177}
\end{align*}

\begin{align*}
\mathcal{Q}_{126}={}&2x_{175}+x_{175}^{2}+x_{176}^{2}-x_{177}^{2}
\end{align*}

\begin{align*}
\mathcal{Q}_{127}={}&-1+x_{116}+x_{138}x_{176}+x_{154}x_{178}-x_{175}\\
&-x_{176}-2x_{176}x_{231}
\end{align*}

\begin{align*}
\mathcal{Q}_{128}={}&-x_{138}x_{180}-x_{180}+x_{181}
\end{align*}

\begin{align*}
\mathcal{Q}_{129}={}&2x_{179}+x_{179}^{2}+x_{180}^{2}-x_{181}^{2}
\end{align*}

\begin{align*}
\mathcal{Q}_{130}={}&-1+x_{117}+x_{138}x_{180}-2x_{138}x_{182}+x_{154}x_{182}\\
&-x_{179}-2x_{180}x_{231}+x_{182}+4x_{182}x_{231}
\end{align*}

\begin{align*}
\mathcal{Q}_{131}={}&-x_{138}x_{184}-x_{184}+x_{185}
\end{align*}

\begin{align*}
\mathcal{Q}_{132}={}&2x_{183}+x_{183}^{2}+x_{184}^{2}-x_{185}^{2}
\end{align*}

\begin{align*}
\mathcal{Q}_{133}={}&-1+x_{120}+x_{138}x_{184}+x_{154}x_{186}-x_{183}\\
&-x_{184}-2x_{184}x_{231}
\end{align*}

\begin{align*}
\mathcal{Q}_{134}={}&-x_{138}x_{188}-x_{188}+x_{189}
\end{align*}

\begin{align*}
\mathcal{Q}_{135}={}&2x_{187}+x_{187}^{2}+x_{188}^{2}-x_{189}^{2}
\end{align*}

\begin{align*}
\mathcal{Q}_{136}={}&-1+x_{123}+x_{138}x_{188}+x_{154}x_{190}-x_{187}\\
&-x_{188}-2x_{188}x_{231}
\end{align*}

\begin{align*}
\mathcal{Q}_{137}={}&1-x_{138}x_{192}-x_{192}+x_{193}
\end{align*}

\begin{align*}
\mathcal{Q}_{138}={}&-1-x_{5}x_{192}+x_{19}+x_{147}x_{148}+x_{191}\\
&+x_{191}^{2}+x_{192}^{2}-x_{193}-x_{193}^{2}
\end{align*}

\begin{align*}
\mathcal{Q}_{139}={}&x_{124}-2x_{138}x_{194}+x_{154}x_{194}-x_{191}-x_{192}\\
&-2x_{192}x_{231}+x_{193}+x_{194}+4x_{194}x_{231}
\end{align*}

\begin{align*}
\mathcal{Q}_{140}={}&-x_{138}x_{196}-x_{196}+x_{197}
\end{align*}

\begin{align*}
\mathcal{Q}_{141}={}&2x_{195}+x_{195}^{2}+x_{196}^{2}-x_{197}^{2}
\end{align*}

\begin{align*}
\mathcal{Q}_{142}={}&-1+x_{128}+x_{138}x_{196}+x_{154}x_{198}-x_{195}\\
&-x_{196}-2x_{196}x_{231}
\end{align*}

\begin{align*}
\mathcal{Q}_{143}={}&-x_{138}x_{200}-x_{200}+x_{201}
\end{align*}

\begin{align*}
\mathcal{Q}_{144}={}&2x_{199}+x_{199}^{2}+x_{200}^{2}-x_{201}^{2}
\end{align*}

\begin{align*}
\mathcal{Q}_{145}={}&-1+x_{131}+x_{138}x_{200}+x_{154}x_{202}-x_{199}\\
&-x_{200}-2x_{200}x_{231}
\end{align*}

\begin{align*}
\mathcal{Q}_{146}={}&-x_{138}x_{204}-x_{204}+x_{205}
\end{align*}

\begin{align*}
\mathcal{Q}_{147}={}&2x_{203}+x_{203}^{2}+x_{204}^{2}-x_{205}^{2}
\end{align*}

\begin{align*}
\mathcal{Q}_{148}={}&-1+x_{134}+x_{138}x_{204}+x_{154}x_{206}-x_{203}\\
&-x_{204}-2x_{204}x_{231}
\end{align*}

\begin{align*}
\mathcal{Q}_{149}={}&-1-2x_{0}-3x_{0}^{2}-2x_{0}x_{4}-x_{1}\\
&-x_{1}^{2}-x_{2}-x_{2}^{2}-x_{4}-2x_{4}^{2}\\
&-x_{5}^{2}-x_{5}x_{20}+x_{5}x_{192}-x_{6}+x_{8}x_{9}\\
&+x_{10}+2x_{10}x_{12}-x_{11}^{2}-x_{12}-x_{13}^{2}\\
&-x_{14}^{2}-x_{15}^{2}-x_{16}^{2}-x_{17}^{2}+x_{18}\\
&-x_{18}^{2}-x_{19}-x_{19}^{2}-x_{26}-x_{98}^{2}\\
&-x_{99}^{2}-x_{103}^{2}-x_{104}^{2}-x_{107}^{2}-x_{110}^{2}\\
&-x_{113}^{2}-x_{116}^{2}-x_{117}^{2}-x_{120}^{2}-x_{123}^{2}\\
&-x_{124}^{2}-x_{128}^{2}-x_{131}^{2}-x_{134}^{2}+x_{137}\\
&-x_{138}+x_{138}^{2}-4x_{138}x_{231}+x_{143}^{2}-x_{147}x_{148}\\
&+2x_{149}x_{150}+x_{152}^{2}+x_{154}-2x_{158}+x_{158}^{2}\\
&-x_{161}+x_{164}^{2}+x_{168}^{2}+x_{172}^{2}+x_{176}^{2}\\
&+x_{180}^{2}+x_{184}^{2}+x_{188}^{2}-x_{191}+x_{192}^{2}\\
&-2x_{192}x_{207}+x_{193}+x_{196}^{2}+x_{200}^{2}-x_{204}^{2}\\
&-x_{207}^{2}+x_{209}x_{215}-x_{213}+x_{225}^{2}+x_{226}^{2}\\
&+x_{227}+x_{227}^{2}-4x_{231}-4x_{231}^{2}
\end{align*}

\begin{align*}
\mathcal{Q}_{150}={}&-x_{5}x_{10}-x_{5}x_{11}-x_{5}x_{12}-x_{5}x_{13}-x_{5}x_{14}\\
&-x_{5}x_{15}-x_{5}x_{16}-x_{5}x_{17}-x_{5}x_{18}+x_{207}
\end{align*}

\begin{align*}
\mathcal{Q}_{151}={}&-2+x_{0}-2x_{1}-x_{2}+2x_{2}x_{7}\\
&-x_{3}-x_{3}x_{10}-x_{3}x_{11}+2x_{3}x_{21}+x_{4}\\
&-x_{5}x_{20}-x_{6}+x_{6}^{2}+x_{11}+x_{11}^{2}\\
&+2x_{11}x_{12}+x_{12}+x_{12}^{2}+x_{18}-x_{25}\\
&-3x_{77}-x_{78}-x_{79}-x_{138}x_{192}-x_{192}\\
&+x_{193}-x_{208}^{2}-x_{209}^{2}-x_{210}^{2}-x_{211}^{2}\\
&-x_{212}^{2}+x_{218}^{2}+x_{219}^{2}
\end{align*}

\begin{align*}
\mathcal{Q}_{152}={}&-1+x_{2}+x_{6}-x_{6}^{2}
\end{align*}

\begin{align*}
\mathcal{Q}_{153}={}&-x_{138}x_{210}-x_{210}+x_{218}
\end{align*}

\begin{align*}
\mathcal{Q}_{154}={}&-1+x_{209}^{2}+x_{210}^{2}-x_{218}^{2}
\end{align*}

\begin{align*}
\mathcal{Q}_{155}={}&-x_{212}-x_{212}x_{213}+x_{219}
\end{align*}

\begin{align*}
\mathcal{Q}_{156}={}&-1+x_{211}^{2}+x_{212}^{2}-x_{219}^{2}
\end{align*}

\begin{align*}
\mathcal{Q}_{157}={}&-x_{204}^{2}+x_{220}
\end{align*}

\begin{align*}
\mathcal{Q}_{158}={}&x_{210}-2x_{214}x_{220}-x_{220}
\end{align*}

\begin{align*}
\mathcal{Q}_{159}={}&x_{138}x_{209}+x_{215}
\end{align*}

\begin{align*}
\mathcal{Q}_{160}={}&-x_{138}-x_{209}x_{215}+x_{213}
\end{align*}

\begin{align*}
\mathcal{Q}_{161}={}&-1-x_{203}-x_{209}x_{216}+x_{211}
\end{align*}

\begin{align*}
\mathcal{Q}_{162}={}&-x_{10}-x_{11}-4x_{21}-2x_{24}-2x_{28}\\
&-2x_{30}-2x_{32}-2x_{34}-4x_{36}-2x_{37}\\
&-2x_{40}-2x_{43}-4x_{62}-2x_{63}-x_{204}x_{217}\\
&+x_{212}
\end{align*}

\begin{align*}
\mathcal{Q}_{163}={}&-1+2x_{0}^{2}+2x_{0}x_{4}-2x_{0}x_{143}-2x_{0}x_{152}\\
&+x_{1}^{2}-2x_{1}x_{192}+x_{2}^{2}-2x_{2}x_{180}+x_{4}^{2}\\
&-2x_{4}x_{152}+x_{10}^{2}-2x_{10}x_{158}-2x_{10}x_{204}+x_{11}^{2}\\
&-2x_{11}x_{164}-2x_{11}x_{204}+x_{12}^{2}-2x_{12}x_{168}+x_{13}^{2}\\
&-2x_{13}x_{172}+x_{14}^{2}-2x_{14}x_{176}+x_{15}^{2}-2x_{15}x_{184}\\
&+x_{16}^{2}-2x_{16}x_{188}+x_{17}^{2}-2x_{17}x_{196}+x_{18}^{2}\\
&-2x_{18}x_{200}-8x_{21}x_{204}-4x_{24}x_{204}-4x_{28}x_{204}-4x_{30}x_{204}\\
&-4x_{32}x_{204}-4x_{34}x_{204}-8x_{36}x_{204}-4x_{37}x_{204}-4x_{40}x_{204}\\
&-4x_{43}x_{204}-8x_{62}x_{204}-4x_{63}x_{204}+x_{138}-x_{138}x_{158}\\
&-x_{138}x_{221}+x_{142}^{2}+x_{151}^{2}+x_{157}^{2}+x_{158}\\
&+x_{159}+x_{163}^{2}+x_{167}^{2}+x_{171}^{2}+x_{175}^{2}\\
&+x_{179}^{2}+x_{183}^{2}+x_{187}^{2}+x_{191}^{2}+x_{195}^{2}\\
&+x_{199}^{2}+x_{203}^{2}+x_{204}^{2}+x_{220}-x_{225}^{2}\\
&-x_{226}^{2}-x_{227}-x_{227}^{2}
\end{align*}

\begin{align*}
\mathcal{Q}_{164}={}&-x_{99}x_{222}+x_{124}
\end{align*}

\begin{align*}
\mathcal{Q}_{165}={}&-x_{104}x_{223}+x_{124}
\end{align*}

\begin{align*}
\mathcal{Q}_{166}={}&-x_{117}x_{224}+x_{124}
\end{align*}

\begin{align*}
\mathcal{Q}_{167}={}&2+2x_{1}-x_{2}x_{7}-x_{3}x_{10}-x_{3}x_{11}\\
&-x_{4}+x_{5}x_{20}-x_{8}x_{9}+x_{10}+x_{10}^{2}\\
&+2x_{10}x_{11}-8x_{10}x_{21}-4x_{10}x_{24}-4x_{10}x_{28}-4x_{10}x_{30}\\
&-4x_{10}x_{32}-4x_{10}x_{34}-8x_{10}x_{36}-4x_{10}x_{37}-4x_{10}x_{40}\\
&-4x_{10}x_{43}-8x_{10}x_{62}-4x_{10}x_{63}+2x_{11}+3x_{11}^{2}\\
&-2x_{11}x_{12}-8x_{11}x_{21}-4x_{11}x_{24}-4x_{11}x_{28}-4x_{11}x_{30}\\
&-4x_{11}x_{32}-4x_{11}x_{34}+4x_{11}x_{36}-4x_{11}x_{37}-4x_{11}x_{40}\\
&-4x_{11}x_{43}-8x_{11}x_{62}-4x_{11}x_{63}-x_{12}-x_{12}^{2}\\
&-x_{18}+8x_{21}-16x_{21}x_{24}-16x_{21}x_{28}-16x_{21}x_{30}\\
&-16x_{21}x_{32}-16x_{21}x_{34}-32x_{21}x_{36}-16x_{21}x_{37}-16x_{21}x_{40}\\
&-16x_{21}x_{43}-32x_{21}x_{62}-16x_{21}x_{63}-2x_{24}-4x_{24}^{2}\\
&-8x_{24}x_{28}-8x_{24}x_{30}-8x_{24}x_{32}-8x_{24}x_{34}-16x_{24}x_{36}\\
&-8x_{24}x_{37}-8x_{24}x_{40}-8x_{24}x_{43}-16x_{24}x_{62}-8x_{24}x_{63}\\
&+x_{25}-2x_{28}-4x_{28}^{2}-8x_{28}x_{30}-8x_{28}x_{32}\\
&-8x_{28}x_{34}-16x_{28}x_{36}-8x_{28}x_{37}-8x_{28}x_{40}-8x_{28}x_{43}\\
&-16x_{28}x_{62}-8x_{28}x_{63}-2x_{30}-4x_{30}^{2}-8x_{30}x_{32}\\
&-8x_{30}x_{34}-16x_{30}x_{36}-8x_{30}x_{37}-8x_{30}x_{40}-8x_{30}x_{43}\\
&-16x_{30}x_{62}-8x_{30}x_{63}-2x_{32}-4x_{32}^{2}-8x_{32}x_{34}\\
&-16x_{32}x_{36}-8x_{32}x_{37}-8x_{32}x_{40}-8x_{32}x_{43}-16x_{32}x_{62}\\
&-8x_{32}x_{63}-2x_{34}-4x_{34}^{2}-16x_{34}x_{36}-8x_{34}x_{37}\\
&-8x_{34}x_{40}-8x_{34}x_{43}-16x_{34}x_{62}-8x_{34}x_{63}-4x_{36}\\
&-4x_{36}^{2}-16x_{36}x_{37}-16x_{36}x_{40}-16x_{36}x_{43}-32x_{36}x_{62}\\
&-16x_{36}x_{63}-5x_{37}-4x_{37}^{2}-8x_{37}x_{40}-8x_{37}x_{43}\\
&-16x_{37}x_{62}-8x_{37}x_{63}-x_{38}-2x_{40}-4x_{40}^{2}\\
&-8x_{40}x_{43}-16x_{40}x_{62}-8x_{40}x_{63}-2x_{43}-4x_{43}^{2}\\
&-16x_{43}x_{62}-8x_{43}x_{63}+4x_{62}-8x_{62}^{2}-16x_{62}x_{63}\\
&-4x_{63}-4x_{63}^{2}+x_{77}+x_{78}+x_{138}x_{192}\\
&-x_{143}^{2}-x_{152}^{2}+2x_{158}-x_{158}^{2}-x_{164}^{2}\\
&-x_{168}^{2}-x_{172}^{2}-x_{176}^{2}-x_{180}^{2}-x_{184}^{2}\\
&-x_{188}^{2}+x_{192}-x_{192}^{2}-x_{193}-x_{196}^{2}\\
&-x_{200}^{2}+x_{204}^{2}-x_{204}x_{217}+x_{212}-x_{225}^{2}\\
&-x_{226}^{2}-x_{227}-x_{227}^{2}
\end{align*}

\begin{align*}
\mathcal{Q}_{168}={}&-1-4x_{101}^{2}-4x_{106}^{2}-4x_{109}^{2}-4x_{112}^{2}\\
&-4x_{115}^{2}-4x_{119}^{2}-4x_{122}^{2}+x_{124}^{2}-4x_{126}^{2}\\
&-4x_{130}^{2}-4x_{133}^{2}-4x_{136}^{2}-x_{228}^{2}-x_{229}^{2}\\
&-x_{230}-x_{230}^{2}-4x_{233}^{2}
\end{align*}

\begin{align*}
\mathcal{Q}_{169}={}&-x_{5}x_{232}+x_{124}
\end{align*}

\begin{align*}
\mathcal{Q}_{170}={}&x_{3}x_{124}-x_{6}x_{232}-x_{233}
\end{align*}

\endgroup
\end{document}
